\section{The integral central-splitting decision}\label{app:central}

This supplies the abelian step of Proposition~\ref{prop:n5}. Given a
finitely generated abelian group $Z$ and finitely many constraints
$v\in Z_i+dZ$, decide whether $Z=Z_1\oplus Z_2$ satisfies them,
including requirements that specified factors be nontrivial.

Write $Z=F\oplus T$, with $F\cong\Z^n$ and $T$ finite. Enumerate
all direct splittings $T=T_1\oplus T_2$ and all homomorphisms
$\beta:F\to T$. These are finite sets. Every ordered splitting of
$Z$ is the image of
\[
 Z_1=A\oplus T_1,\qquad Z_2=B\oplus T_2,\qquad F=A\oplus B
\]
under a shear $(f,t)\mapsto(f,t+\beta(f))$. Indeed the intersections
with $T$ split $T$, the images in $F$ split $F$, and chosen free
summands are graphs of homomorphisms into the complementary torsion
factor. Combining these graph maps gives $\beta$.

Apply the inverse shear to the constraints. Each separates into the
finite test $t\in T_i+dT$ and a free test $f\in A_i+dF$.
Let $W_i$ be the saturation of the lattice spanned by the exact
($d=0$) vectors assigned to side $i$. The exact conditions are possible
precisely when
\[
 W_1\cap W_2=0,\qquad W_1+W_2\text{ is primitive in }F.
\]
Test these by Smith form and choose $F=W_1\oplus W_2\oplus C$.
Every splitting containing the $W_i$ is obtained from one of
\[
 A_k=W_1\oplus\langle c_1,\ldots,c_k\rangle,\qquad
 B_k=W_2\oplus\langle c_{k+1},\ldots,c_e\rangle,\quad0\le k\le e,
\]
where $e=\rank C$, by an automorphism fixing $W_1+W_2$ pointwise.
Such automorphisms have matrices
\[
 \Gamma=\left\{
 \begin{pmatrix}I&0&U\\0&I&V\\0&0&W\end{pmatrix}:
 W\in\GL_e(\Z)\right\}.
\]
A finite generating set consists of the displayed shears, elementary
matrices, signed permutations on $C$. Let $M$ be the lcm of all
positive moduli, or one if there are none. Enumerate the finite
image of $\Gamma$ modulo $M$, retaining an actual word in integral
generators for each image. For each image and each $k$, compute the
projection $P$ onto the transported $A_k$ along $B_k$. The remaining
constraints are exactly
\[
 (1-P)f=0\pmod d\quad\text{on side one},\qquad
 Pf=0\pmod d\quad\text{on side two}.
\]
They are finite congruence tests. A passing stored word gives an actual
integral splitting, not merely a modular idempotent; the free ranks and
torsion sizes decide required nontriviality. Conversely every splitting
is represented by one of these branches and residues. This proves the
finite constructive step used in the group argument. In particular,
idempotents with incompatible ranks at different primes cannot be
mistaken for an integral splitting.

For the preceding Smith reduction, if $UE_iV=D_i$, apply $U$ to the
central defect vector. A zero row requires membership in $Z_i$;
a diagonal entry $d>0$ requires membership in $Z_i+dZ$. These operations
are valid over every abelian group, including finite groups. Once a
splitting passes, solving the same Smith system constructs corrected
lifts and therefore the group factors.

The controlling \record{problems/N5/proof.md}{proof} and
\record{problems/N5/integrality-audit.md}{integrality audit} include
torsion shears, mixed centres, composite moduli and dependent-relator
controls. The later
\record{problems/N5/general-fp-audit.md}{finite-presentation command}
connects this stage to rational support and the complete quotient search.
